\documentclass[11pt]{article}
\usepackage{amsmath,amssymb,amsthm}
\usepackage[T1]{fontenc}
\usepackage[utf8]{inputenc}
\usepackage{amsmath,amssymb,amsthm,mathtools}
\usepackage[margin=1in]{geometry}
\usepackage{microtype}
\usepackage[hidelinks]{hyperref}

\newtheorem{theorem}{Theorem}[section]
\newtheorem{proposition}[theorem]{Proposition}
\newtheorem{lemma}[theorem]{Lemma}
\theoremstyle{remark}
\newtheorem{remark}[theorem]{Remark}

\newcommand{\R}{\mathbb R}
\newcommand{\E}{\mathbb E}
\newcommand{\Prob}{\mathbb P}
\newcommand{\one}{\mathbf 1}
\newcommand{\cF}{\mathcal F}
\newcommand{\Sph}{\mathbb S^{d-1}}
\DeclareMathOperator{\divergence}{div}
\DeclareMathOperator{\KL}{KL}
\DeclareMathOperator{\Id}{Id}

\newcommand{\fptitle}{Review of Problem 2 -- Astra}
\newcommand{\reviewtag}{Reviewer 2}
\usepackage{fp-review}

\begin{document}
\maketitle

\section{Review}

\noindent \textbf{Summary.} The manuscript combines the inversion formula for the Darcy problem with a plug-in kernel estimator for the solution field $u_f$ to establish an upper bound on the estimation risk. This is accomplished without a pre-estimation step. On a technical level, the analysis combines properties of kernel estimators with stability properties of the inverse Darcy operator. The lower bound is established using a classical lower-bound strategy based on perturbations of the solution field $u_f$. The upper and lower bounds match up to logarithmic factors.

\noindent \textbf{Novelty.} The general approach to establishing the upper bounds relies on the inversion formula for the Darcy problem combined, as in the human solution. The use of this formula is itself novel in the field. The main difference from the human solution lies in the direct treatment of the $L^2$-risk, rather than the pointwise risk, and in a different error decomposition. There are also minor differences in the construction of the estimator.

\paragraph{Correctness.}
\begin{itemize}
	\item  1: essentially flawless mathematical argument
	\item  0.5: minor errors that are easily corrected
	\item  0: blank solution, or acknowledged partial result
	\item  $-0.5$: major errors
	\item  $-1$: proposed proof makes no progress towards the result.
\end{itemize}
Note that points should be docked for wrong statements/ solutions.

\noindent Score: 0.5

\noindent Justification: The manuscript establishes the correct upper and lower bounds (up to logarithmic factors). While the general estimation and proof methodology is sound and correctly laid out, the necessary details of the proofs are missing. A concrete example is provided by the proof of Lemma 4.3, which is more of a sketch than a rigorous mathematical argument.

\paragraph{Exposition.}
\begin{itemize}
	\item  1: well-written and clear solution.
	\item  0: argument that is possible to follow with serious effort
	\item  $-1$: very poor exposition
\end{itemize}

\noindent Score: 0

\noindent Justification: While the paper provides valuable mathematical ideas, it lacks sufficient guidance for the reader. For example, the PDE model is missing. Moreover, the underlying reasons for the non-standard convergence rate are only insufficiently explained in Remark 8.3, at the very end of the manuscript. Some proofs would also benefit from further explanation, in addition to the already mentioned need for more mathematical detail.

\paragraph{Attributions.}
\begin{itemize}
	\item  1: solutions contains fully appropriate attributions
		to related work and theorems being used.
	\item  0: moderate issues with missing attributions or hallucinated references
	\item  $-1$: blatant plagiarism or hallucinations
\end{itemize}

\noindent Score: 0

\noindent Justification: The paper lacks a literature review. Of the four references provided, only one is statistical in nature. Some results are reproved without proper attribution (e.g., the inversion formulas and properties of kernel estimators).

\newpage
\section{Submission with inline comments}


\begin{abstract}
For a sufficiently small, fixed H\"older--Zygmund ball of conductivities
around the constant function \(1\), the minimax squared \(L^2\)-risk is,
up to logarithmic factors,
\[
n^{-2\alpha/(2\alpha+d+3)}.
\]
We prove matching polynomial upper and lower bounds for every \(d\geq2\)
and \(\alpha>1\). The upper bound uses a regularized reconstruction along
gradient trajectories, together with a variance estimate that exploits
integration along these trajectories. The lower bound uses perturbations
that oscillate in the \(d-1\) angular directions. No boundary condition
on the conductivity is imposed.
\end{abstract}

\section{Statement and conventions}

Let \(D=B(0,1)\subset\R^d\). We interpret \(C^s(D)\) as the
H\"older--Zygmund restriction space \(B^s_{\infty,\infty}(D)\).
Any fixed usual equivalent norm may be used; the constants below,
including the admissible radius, may depend on this choice.
\review{The introduction of the PDE model is missing.}
White noise means the isonormal experiment
\[
Y_n(\varphi)=\langle u_f,\varphi\rangle_{L^2(D)}
+n^{-1/2}\mathbb W(\varphi),\qquad \varphi\in L^2(D),
\]
where
\(\E[\mathbb W(\varphi)\mathbb W(\psi)]=\langle\varphi,\psi\rangle\).
\review{Inconsistent notation of inner products.}
Write
\[
\cF(\alpha,r)=
\{f:\|f-\one\|_{C^\alpha(D)}\leq r\},
\qquad
R_n(\alpha,r)=
\inf_{\widehat f}\sup_{f\in\cF(\alpha,r)}
\E_f\|\widehat f-f\|_{L^2(D)}^2.
\]

\begin{theorem}\label{thm:main}
Fix \(d\geq2\) and \(\alpha>1\). There exists \(r_*>0\) such that,
for each fixed \(r\in(0,r_*]\), there are constants
\(c,C>0\) and \(n_0\) for which
\[
c\,n^{-2\alpha/(2\alpha+d+3)}
\leq R_n(\alpha,r)
\leq
C(\log n)^4 n^{-2\alpha/(2\alpha+d+3)},
\qquad n\geq n_0.
\]
If \(\alpha\notin\mathbb N\), the logarithmic factor in this upper bound
can be omitted. In particular, the answer up to logarithmic factors is
\[
\boxed{\displaystyle
R_n(\alpha,r)=\widetilde{\asymp}\,
n^{-2\alpha/(2\alpha+d+3)}.}
\]
\end{theorem}
\review{Notation $=\widetilde{\asymp}$ faulty. Notations for $\partial^\beta$ and $\Delta_j$ (below) not properly defined. The role of $a$ changes throughout the manuscript, often without defining its domain.}
For a smooth function \(a\), the notation
\[
\|a\|_{k,\infty;\Omega}
=\max_{|\beta|\leq k}\|\partial^\beta a\|_{L^\infty(\Omega)}
\]
denotes the classical derivative norm, not an integer-index Zygmund
norm. Domain subscripts will be suppressed when unambiguous.
Constants denoted by \(C\) may change from line to line, but are
independent of \(n\), the bandwidth, and the conductivity in the
specified parameter class.

\review{Literature missing.}

\section{Analytic preliminaries}

\subsection{Function spaces and the forward equation}

We use the Littlewood--Paley characterization
\begin{equation}\label{eq:LP}
\|a\|_{C^s(\R^d)}
\asymp \sup_{j\geq0}2^{js}\|\Delta_j a\|_\infty
\end{equation}
for an inhomogeneous smooth dyadic decomposition. The equivalence with
classical H\"older spaces at noninteger positive indices, and the
differentiation properties of these spaces, are given in
\cite[Chapter~13, Section~8]{Taylor}.

We also use a common bounded extension operator from a smooth bounded
domain to \(\R^d\), bounded on all the finitely many Besov spaces under
consideration. Such an operator exists even for Lipschitz domains;
see \cite{Rychkov}.

Here are some consequences, with the relevant arguments recalled.
Differentiation maps \(C^s\) boundedly to \(C^{s-1}\).
For \(s>0\), \(C^s\) is a multiplication algebra:
\begin{equation}\label{eq:algebra}
\|ab\|_{C^s}\leq C_s\|a\|_{C^s}\|b\|_{C^s}.
\end{equation}
Indeed, decompose the product into low--high, high--low, and
comparable-frequency products. The low--high contribution to its
\(j\)-th dyadic block is bounded by
\(C\|a\|_\infty 2^{-js}\|b\|_{C^s}\); interchange \(a,b\) for
the high--low contribution. For comparable frequencies \(k\geq j-C\),
the bound is
\[
C\|a\|_{C^s}\|b\|_{C^s}
\sum_{k\geq j-C}2^{-2ks}
\leq C_s2^{-js}\|a\|_{C^s}\|b\|_{C^s}.
\]
This proves \eqref{eq:algebra}. Extension and restriction give the same
facts on \(D\).

We shall also use the real interpolation identity
\begin{equation}\label{eq:interp}
(C^{s_0},C^{s_1})_{\theta,\infty}
=C^{(1-\theta)s_0+\theta s_1},
\qquad s_0<s_1,\quad 0<\theta<1.
\end{equation}
For completeness, on \(\R^d\) this follows by splitting the dyadic
series at a level \(J\). If \(s=(1-\theta)s_0+\theta s_1\), the low
and high parts satisfy
\[
\|a_{\leq J}\|_{C^{s_1}}
\leq C2^{J(s_1-s)}\|a\|_{C^s},\qquad
\|a_{>J}\|_{C^{s_0}}
\leq C2^{-J(s-s_0)}\|a\|_{C^s}.
\]
Taking \(2^{J(s_1-s_0)}\asymp t^{-1}\) bounds the \(K\)-functional
by \(Ct^\theta\|a\|_{C^s}\). Conversely, applying \(\Delta_j\) to
any decomposition \(a=a_0+a_1\) gives
\[
\|\Delta_j a\|_\infty
\leq C2^{-js_0}
K(2^{-j(s_1-s_0)},a;C^{s_0},C^{s_1}),
\]
which proves the reverse bound. A common extension operator transfers
the identity to \(D\).

\begin{lemma}[Uniform perturbative regularity]\label{lem:forward}
Put
\[
u_\circ(x)=\frac{|x|^2-1}{2d}.
\]
For sufficiently small \(r_*>0\), every
\(f\in\cF(\alpha,r_*)\) satisfies \(1/2\leq f\leq3/2\), and the
Dirichlet problem
\[
\divergence(f\nabla u)=1,\qquad u|_{\partial D}=0
\]
has a unique weak solution. Moreover,
\begin{equation}\label{eq:forward-reg}
\|u_f-u_\circ\|_{C^{\alpha+1}(D)}
\leq C\|f-\one\|_{C^\alpha(D)}.
\end{equation}
\end{lemma}

\begin{proof}
Let \(\mathcal P F\) be the zero-Dirichlet solution of
\(\Delta w=\divergence F\). For every \(s>1\),
\begin{equation}\label{eq:poisson}
\|\mathcal P F\|_{C^{s+1}(D)}
\leq C_s\|F\|_{C^s(D)}.
\end{equation}
At noninteger \(s\), this is the classical global Schauder estimate
and solvability theorem, including higher-order regularity; see
\cite[Theorems~6.6, 6.14 and 6.19]{GT}.
At integer \(s\), apply these bounds at two surrounding noninteger
indices and use \eqref{eq:interp}. Thus \eqref{eq:poisson} holds
also at the Zygmund indices.

Writing \(u=u_\circ+w\), the desired equation becomes
\[
w=-\mathcal P\big((f-1)\nabla u_\circ\big)
  -\mathcal P\big((f-1)\nabla w\big).
\]
By \eqref{eq:algebra} and \eqref{eq:poisson}, the second term is a
linear operator on the zero-trace subspace of \(C^{\alpha+1}(D)\)
with norm at most \(C\|f-1\|_{C^\alpha}\). Choosing \(r_*\) so that
this norm is at most \(1/2\), the Neumann series gives a solution
and \eqref{eq:forward-reg}.

The embedding \(C^\alpha(D)\hookrightarrow C^0(\overline D)\)
allows us to decrease \(r_*\) so that \(1/2\leq f\leq3/2\).
Uniqueness of a weak solution follows by testing the equation for
the difference of two solutions against that difference:
\(\int_D f|\nabla w|^2=0\). Finally, the constructed solution has
zero trace and belongs to \(H^1_0(D)\), so it is that unique weak
solution.
\end{proof}

\subsection{Smoothing and a product cancellation}

Set
\[
m_0=\lceil\alpha\rceil-1.
\]
Fix a real \(K\in C^\infty_c(B(0,1))\) such that
\begin{equation}\label{eq:kernel}
\int K=1,\qquad
\int x^\beta K(x)\,dx=0
\quad(1\leq|\beta|\leq2m_0+2).
\end{equation}
Such a kernel can be obtained from a finite linear combination of
distinct dilations of a smooth unit-integral bump: the required
coefficients solve a Vandermonde system in the dilation parameters.
Positivity of \(K\) is neither required nor asserted.

For \(0<h<1/8\), put
\[
K_h(x)=h^{-d}K(x/h),\qquad
\Omega_h=B(0,1-2h),\qquad
S_h a(x)=\int_D K_h(x-y)a(y)\,dy,\quad x\in\Omega_h,
\]
and let \(L_h=\log(e/h)\).

\begin{lemma}[Smoothing and commutator bounds]\label{lem:commutator}
Suppose \(a,b\in C^\alpha(D)\) have bounded norms. Uniformly on
\(\Omega_h\),
\begin{align}
\|S_h a-a\|_\infty&\leq C h^\alpha L_h,
\label{eq:smoothing}\\
\left\|\partial^\gamma
\big((S_h a)(S_h b)-S_h(ab)\big)\right\|_\infty
&\leq C h^{\alpha+1-|\gamma|}L_h^2,
\qquad |\gamma|\leq2.
\label{eq:commutator}
\end{align}
If \(\alpha\notin\mathbb N\), all factors \(L_h\) in these two
bounds may be replaced by \(1\). The statements hold componentwise
for vector-valued \(b\).
\end{lemma}

\begin{proof}
We first record the local Taylor remainder
\begin{equation}\label{eq:taylor}
\sup_{|y-x|\leq h}|a(y)-P_{a,x}(y)|
\leq C h^\alpha L_h,
\end{equation}
where \(P_{a,x}\) is the Taylor polynomial of degree \(m_0\).
Its coefficients are uniformly bounded. At noninteger \(\alpha\),
this is the ordinary H\"older Taylor estimate, without \(L_h\).
At integer \(\alpha\), derivatives of order \(\alpha-1\) have
modulus of continuity bounded by \(Ct\log(e/t)\). To see the latter
directly from \eqref{eq:LP}, split their dyadic series at
\(2^J\asymp t^{-1}\): the low-frequency increments sum to
\(C\sum_{j\leq J}t\), and the high-frequency tail is at most
\(C\sum_{j>J}2^{-j}\). Taylor's integral remainder then gives
\eqref{eq:taylor}. Extension to \(\R^d\) justifies this argument
uniformly near the boundary.

Polynomial reproduction from \eqref{eq:kernel} immediately proves
\eqref{eq:smoothing}. To prove \eqref{eq:commutator}, write
\[
Q_h(a,b)=(S_h a)(S_h b)-S_h(ab).
\]
Fix \(x\) and write \(a=P_a+R_a\), \(b=P_b+R_b\), using the Taylor
polynomials at \(x\). Keep these polynomials fixed when differentiating
at \(x\). Polynomial reproduction gives \(Q_h(P_a,P_b)=0\).
Also \(Q_h(a,c)=0\) for every constant \(c\).

On \(B(x,h)\), the remainders satisfy
\(|R_a|+|R_b|\leq Ch^\alpha L_h\), whereas
\(|P_a-P_a(x)|+|P_b-P_b(x)|\leq Ch\).
Differentiating the convolution kernels gives, for \(|\gamma|\leq2\),
\[
|\partial^\gamma Q_h(R_a,P_b)(x)|
+
|\partial^\gamma Q_h(P_a,R_b)(x)|
\leq C h^{\alpha+1-|\gamma|}L_h.
\]
For example, in the first expression subtract the constant \(P_b(x)\);
the differentiated convolution of
\(R_a(P_b-P_b(x))\) has precisely this bound. In the product term,
a derivative falling on the polynomial is bounded by
\(Ch^{1-k}\) at derivative order \(k\geq1\); the undifferentiated
polynomial vanishes at \(x\). Finally,
\[
|\partial^\gamma Q_h(R_a,R_b)(x)|
\leq C h^{2\alpha-|\gamma|}L_h^2.
\]
Since \(\alpha>1\), these inequalities prove
\eqref{eq:commutator}. The same proof without logarithms applies
at noninteger \(\alpha\).
\end{proof}

\begin{lemma}[Completion of a thin boundary strip]\label{lem:extension}
For all sufficiently small \(h\), there is a linear map
\(E_h:L^2(\Omega_h)\to L^2(D)\) such that
\begin{align}
\|E_h a\|_{L^2(D)}&\leq C\|a\|_{L^2(\Omega_h)},\label{eq:E-stable}\\
\|E_h(f|_{\Omega_h})-f\|_{L^2(D)}
&\leq C\|f\|_{C^\alpha(D)}h^\alpha L_h.
\label{eq:E-approx}
\end{align}
At noninteger \(\alpha\), \(L_h\) can be omitted.
\end{lemma}

\begin{proof}
Let \(q=m_0+1\) and
\(\lambda_j=(-1)^{j+1}\binom qj\), \(1\leq j\leq q\).
Then
\[
\sum_{j=1}^q\lambda_j=1,\qquad
\sum_{j=1}^q\lambda_j j^k=0,\quad 1\leq k\leq m_0.
\]
Define \(E_h a=a\) on \(\Omega_h\), and, for \(x\in D\setminus\Omega_h\),
define
\[
E_h a(x)=\sum_{j=1}^q\lambda_j a((1-3jh)x).
\]
For \(h<1/(6q)\), all arguments on the right belong to \(\Omega_h\)
and all dilation factors are at least \(1/2\).
A change of variables proves \eqref{eq:E-stable}.
Taylor expansion of \(f\) at \(x\), together with the displayed
moment identities and \eqref{eq:taylor}, proves
\eqref{eq:E-approx}.
\end{proof}

\section{Inversion along gradient trajectories}

For a smooth function \(v\), define
\[
T_v a=\divergence(a\nabla v)
=\nabla v\cdot\nabla a+(\Delta v)a.
\]

\begin{lemma}[Transport inverse]\label{lem:transport}
Let \(\Omega=B(0,R)\), \(3/4\leq R\leq1\), and suppose \(v\) is
smooth on a neighborhood of \(\overline\Omega\), with
\begin{equation}\label{eq:geometry}
\lambda I\leq D^2v\leq\Lambda I,\qquad
\nabla v\cdot\nu>0\ \text{on }\partial\Omega,
\end{equation}
where \(0<\lambda\leq\Lambda\).
There is a unique critical point \(p\in\Omega\).
Let \(z_t(x)\) solve
\[
\frac{d}{dt}z_t(x)=-\nabla v(z_t(x)),\qquad z_0(x)=x.
\]
This flow exists in \(\overline\Omega\) for all \(t\geq0\), and
\begin{equation}\label{eq:radial-flow}
e^{-\Lambda t}|x-p|
\leq |z_t(x)-p|
\leq e^{-\lambda t}|x-p|.
\end{equation}

For smooth \(g\), the unique bounded smooth solution of \(T_v a=g\)
is
\begin{equation}\label{eq:transport-inverse}
P_vg(x)=
\int_0^\infty
\exp\left\{-\int_0^t\Delta v(z_s(x))\,ds\right\}
g(z_t(x))\,dt.
\end{equation}
With \(c_0=d\lambda\), it satisfies
\begin{align}
\|P_vg\|_\infty&\leq c_0^{-1}\|g\|_\infty,\label{eq:P-infty}\\
\|P_vg\|_2&\leq 2c_0^{-1}\|g\|_2,\label{eq:P-L2}\\
\|P_vg\|_{H^1}
&\leq C\big\{\|g\|_{H^1}
 +(1+\|D^3v\|_\infty)\|g\|_2\big\}.
\label{eq:P-H1}
\end{align}
The constants in the last estimate depend only on \(d,\lambda,\Lambda\).
\end{lemma}
\review{Attribution for the inversion formula missing.}

\begin{proof}
Strict convexity gives at most one critical point. A minimum on
\(\overline\Omega\) cannot occur at the boundary, because the inward
direction decreases \(v\), so this critical point exists in \(\Omega\).
The vector field \(-\nabla v\) points inward at the boundary;
the usual first-exit argument \review{Details or reference required.} and local existence for a smooth
vector field therefore give a global backward-gradient flow in
\(\overline\Omega\).

For \(z\in\Omega\),
\[
\lambda|z-p|^2
\leq (z-p)\cdot\nabla v(z)
\leq\Lambda|z-p|^2.
\]
Applying this to the derivative of \(|z_t-p|\) proves
\eqref{eq:radial-flow}.

Since \(\Delta v\geq c_0\), the integral in
\eqref{eq:transport-inverse} converges absolutely and gives
\eqref{eq:P-infty}. It is smooth up to the boundary. Here is a
justification of differentiating the infinite-time integral.
The first variational equation gives
\(\|D_xz_t\|\leq e^{-\lambda t}\).
Inductively, the equation for \(D_x^kz_t\) has homogeneous part
\(-D^2v(z_t)D_x^kz_t\), while each forcing term is a product of
at least two lower-order flow derivatives. Variation of constants
therefore gives
\(\|D_x^kz_t\|\leq C_ke^{-\lambda t}\).
Derivatives of \(\int_0^t\Delta v(z_s)\,ds\) of every positive
order are consequently bounded uniformly in \(t\).
Every differentiated integrand is thus bounded by \(C_ke^{-c_0t}\).
\review{Argument is lacking mathematical precision and more details should be given. This issue persists throughout the manuscript.}

The flow semigroup property, or differentiation after splitting the
integral at a finite time, gives \(T_vP_vg=g\).
If \(T_va=0\), differentiation along the flow shows
\[
a(x)=
\exp\left\{-\int_0^t\Delta v(z_s(x))\,ds\right\}a(z_t(x)).
\]
Letting \(t\to\infty\) proves uniqueness among bounded solutions.

For \(a=P_vg\), integration by parts gives
\[
\int_\Omega aT_va
=
\frac12\int_{\partial\Omega}(\nabla v\cdot\nu)a^2
+\frac12\int_\Omega(\Delta v)a^2
\geq\frac{c_0}{2}\|a\|_2^2.
\]
This proves \eqref{eq:P-L2}. Differentiating the equation gives
\[
T_v(\nabla a)+D^2v\,\nabla a
=\nabla g-a\nabla\Delta v.
\]
The same integration by parts, now for a vector field, yields
\[
(c_0/2+\lambda)\|\nabla a\|_2
\leq\|\nabla g\|_2+\|\nabla\Delta v\|_\infty\|a\|_2.
\]
Together with \eqref{eq:P-L2}, this proves \eqref{eq:P-H1}.
\end{proof}

\section{Gaussian smoothing and the decisive variance estimate}

Define
\[
Z_h(x)=n^{-1/2}\mathbb W(K_h(x-\cdot)),\qquad x\in\R^d,
\]
where the white-noise argument is restricted to \(D\).

\begin{lemma}[Gaussian smoothing bounds]\label{lem:gaussian}
For every fixed nonnegative integer \(k\) and \(p\geq1\),
\begin{align}
\E\|Z_h\|_{k,\infty;\Omega_h}^{p}
&\leq C_p n^{-p/2}h^{-p(d/2+k)}L_h^{p/2},
\label{eq:gaussian-sup}\\
\E\|Z_h\|_{H^k(\Omega_h)}^2
&\leq Cn^{-1}h^{-d-2k},\label{eq:gaussian-H2}\\
\E\|Z_h\|_{H^k(\Omega_h)}^4
&\leq Cn^{-2}h^{-2d-4k}.
\label{eq:gaussian-H4}
\end{align}
Moreover, for \(t\geq0\),
\begin{equation}\label{eq:gaussian-tail}
\Prob\left(
\|Z_h\|_{k,\infty;\Omega_h}>
C n^{-1/2}h^{-d/2-k}\sqrt{L_h+t}
\right)\leq Ce^{-t}.
\end{equation}
\end{lemma}

\begin{proof}
The white-noise isometry gives
\[
\E|\partial^\beta Z_h(x)|^2
\leq Cn^{-1}h^{-d-2|\beta|}.
\]
Integrating proves \eqref{eq:gaussian-H2}.
For centered jointly Gaussian variables \(X,Y\),
\(\E(X^2Y^2)\leq3(\E X^2)(\E Y^2)\);
double integration proves \eqref{eq:gaussian-H4}.

The field has a smooth version. Indeed, after multiplication by a
fixed smooth cutoff, its weak derivatives of every order have finite
expected squared \(L^2\)-norm. Almost surely it therefore belongs to
every integer Sobolev space locally. Fourier inversion and
Cauchy--Schwarz imply \(H^s(\R^d)\hookrightarrow C^k\) for
\(s>k+d/2\), which gives the assertion.

For a derivative of order \(k\), let
\(A=n^{-1/2}h^{-d/2-k}\).
Its Gaussian increment metric is bounded by
\[
CA\min\{1,|x-y|/h\}.
\]
On a fixed cube containing \(\Omega_h\), use successive grids of mesh
\(h2^{-j}\). The initial grid has at most \(Ch^{-d}\) points, and
the \(j\)-th grid has at most \(Ch^{-d}2^{jd}\) points.
Join each point to a nearest point of the preceding grid.
Gaussian tail bounds and a union bound show that, outside an event
of probability \(Ce^{-t}\), the initial values are bounded by
\(CA\sqrt{L_h+t}\), and every increment at level \(j\geq1\) is bounded by
\[
CA2^{-j}\sqrt{L_h+t+j}.
\]
The constants can be chosen so that the failure probabilities at
level \(j\) sum to \(Ce^{-t}\).
Summing the increments and using continuity proves
\eqref{eq:gaussian-tail}. A finite union handles all derivatives
of order at most \(k\). Integrating the tail proves
\eqref{eq:gaussian-sup}.
\end{proof}

The next lemma is responsible for the denominator \(2\alpha+d+3\).

\begin{lemma}[Variance after transport inversion]\label{lem:variance}
Suppose \(v\) satisfies \eqref{eq:geometry} on \(\Omega_h\), and
\[
2d\lambda>\Lambda.
\]
If \(\|a\|_{1,\infty;\Omega_h}\leq A_0\), then
\begin{equation}\label{eq:key-variance}
\E\left\|P_v\divergence(a\nabla Z_h)\right\|_{L^2(\Omega_h)}^2
\leq Cn^{-1}h^{-d-3}.
\end{equation}
The constant depends only on \(d,\lambda,\Lambda,A_0,K\), and not
on higher derivatives of \(v\) or \(a\).
\end{lemma}

\begin{proof}
For \(z\in\Omega_h\), set
\[
H_h(z,y)=a(z)\Delta K_h(z-y)+\nabla a(z)\cdot\nabla K_h(z-y).
\]
Its \(L^2(D)\)-norm is at most \(Ch^{-d/2-2}\), and its support
in \(y\) is contained in \(B(z,h)\). Consequently,
\begin{equation}\label{eq:kernel-overlap}
|\langle H_h(z,\cdot),H_h(w,\cdot)\rangle|
\leq Ch^{-d-4}\one_{\{|z-w|\leq2h\}}.
\end{equation}
Let \(p\) be the critical point of \(v\), \(c_0=d\lambda\), and
\(r=|x-p|\). The representation \eqref{eq:transport-inverse},
the white-noise isometry, and \eqref{eq:kernel-overlap} give the
pointwise variance bound
\[
Cn^{-1}h^{-d-4}
\int_0^\infty\int_0^\infty
e^{-c_0(t+s)}
\one_{\{|z_t(x)-z_s(x)|\leq2h\}}\,ds\,dt.
\]

For \(s=t+w\), put \(r_t=|z_t(x)-p|\).
By \eqref{eq:radial-flow},
\[
|z_t-z_{t+w}|
\geq r_t-r_{t+w}
\geq (1-e^{-\lambda w})r_t.
\]
If \(r_t>4h\), overlap therefore forces
\(w\leq4h/(\lambda r_t)\). If \(r_t\leq4h\), the exponential
integral in \(w\) is bounded by a constant. Thus
\[
\int_0^\infty e^{-c_0w}
\one_{\{|z_t-z_{t+w}|\leq2h\}}\,dw
\leq C\min\{1,h/r_t\}.
\]
Since \(r_t\geq e^{-\Lambda t}r\), the double time integral is at most
\[
C\int_0^\infty e^{-2c_0t}
\min\{1,(h/r)e^{\Lambda t}\}\,dt
\leq C\min\{1,h/r\}.
\]
The last inequality uses \(2c_0>\Lambda\); at \(r=0\) the minimum
is interpreted as \(1\).

Finally, since \(d\geq2\) and \(\Omega_h\subset B(p,2)\),
\[
\int_{\Omega_h}\min\{1,h/|x-p|\}\,dx
\leq C\left(h^d+h\int_h^2 r^{d-2}\,dr\right)
\leq Ch.
\]
Integrating the pointwise variance proves \eqref{eq:key-variance}.
All interchanges of integrals are justified by the exponential
weights and the uniform \(L^2\)-bound on \(H_h\).
\end{proof}

\section{The upper bound}

\begin{proposition}\label{prop:upper}
For sufficiently small \(r_*>0\), the upper bound in
Theorem~\ref{thm:main} holds uniformly over
\(\cF(\alpha,r_*)\).
\end{proposition}

\begin{proof}
Write
\[
p_\alpha=2\alpha+d+3,\qquad h=n^{-1/p_\alpha}.
\]
Only sufficiently large \(n\) will be considered.

\subsection*{Construction of the estimator}

Form the observable smooth field
\[
\widehat u_h(x)=Y_n(K_h(x-\cdot)),\qquad x\in\Omega_h.
\]
For a fixed true \(f\), use the deterministic notation
\[
u=u_f,\qquad v=S_hu,\qquad a=S_hf.
\]
Then \(\widehat u_h=v+Z_h\).

Choose a fixed \(\delta>0\), depending only on \(d\), so small that
\[
\|w-u_\circ\|_{2,\infty;\Omega_h}\leq\delta
\]
implies \eqref{eq:geometry}, uniformly for \(h<1/8\), with
\[
\lambda=\frac{3}{4d},\qquad
\Lambda=\frac{5}{4d}.
\]
Indeed, sufficiently small entrywise perturbations of \(D^2u_\circ=I/d\)
give these Hessian bounds, and
\(\nabla u_\circ\cdot\nu=(1-2h)/d\) on \(\partial\Omega_h\).
These constants also satisfy \(2d\lambda>\Lambda\).

Polynomial reproduction and Lemma~\ref{lem:forward} imply
\[
\|v-u_\circ\|_{2,\infty}\leq Cr_*.
\]
Also, by moving one derivative onto the kernel,
\[
\|a\|_{1,\infty}\leq C,\qquad
\|a\|_{2,\infty}\leq Ch^{-1},\qquad
\|v\|_{3,\infty}\leq Ch^{-1}.
\]
Decrease \(r_*\) so that the first bound is at most \(\delta/2\).
Choose a fixed sufficiently large \(M_0\), and define the observable
event
\[
\mathcal A_h=
\left\{
\|\widehat u_h-u_\circ\|_{2,\infty}\leq\delta,\quad
h\|\widehat u_h\|_{3,\infty}\leq M_0
\right\}.
\]
Lemma~\ref{lem:gaussian} implies, uniformly in the parameter class,
\begin{equation}\label{eq:good-event}
\Prob_f(\mathcal A_h^c)\leq Cn^{-3}.
\end{equation}
To verify this explicitly, take \(t=3\log n\) in
\eqref{eq:gaussian-tail}. Both the \(C^2\)-noise norm and \(h\)
times the \(C^3\)-noise norm are bounded, with this probability,
by a constant times
\[
n^{-1/2}h^{-d/2-2}\sqrt{\log n}
=h^{\alpha-1/2}\sqrt{\log n}=o(1).
\]

On \(\mathcal A_h\), define
\[
\widetilde f_h=P_{\widehat u_h}\one
\quad\text{on }\Omega_h.
\]
Let \(\Pi\) denote pointwise projection onto \([1/2,3/2]\), \review{The reason for the projection should be given.} and set
\[
\widehat f_n=
\begin{cases}
\Pi(E_h\widetilde f_h),&\mathcal A_h,\\
\one,&\mathcal A_h^c.
\end{cases}
\]
This is an estimator based only on the observations.
Its measurability follows from the flow formula:
finite-time solutions of the smooth ODE depend continuously on the
vector field, and the tail of \eqref{eq:transport-inverse} is uniformly
exponentially small on \(\mathcal A_h\).
Notice that the auxiliary transport reconstruction does not require
a boundary value for \(\widetilde f_h\).

\review{The following paragraphs discuss the error decomposition, which should be given here for the reader's convenience.}

\subsection*{Deterministic smoothing bias}

Set
\[
\mathcal R_h=1-T_va.
\]
Because \(\divergence(f\nabla u)=1\), convolution commutes with
differentiation in the interior and gives
\[
\mathcal R_h
=-\divergence\left((S_hf)(S_h\nabla u)-S_h(f\nabla u)\right).
\]
Both \(f\) and \(\nabla u\) have uniformly bounded \(C^\alpha\)-norms.
Lemma~\ref{lem:commutator} therefore gives, with
\[
B_h=h^\alpha L_h^2,
\]
the bounds
\begin{equation}\label{eq:residual-bounds}
\|\mathcal R_h\|_2\leq CB_h,\qquad
\|\mathcal R_h\|_{H^1}\leq Ch^{-1}B_h.
\end{equation}
Also, Lemmas~\ref{lem:commutator} and \ref{lem:extension} give
\begin{equation}\label{eq:completed-bias}
\|E_ha-f\|_{L^2(D)}\leq Ch^\alpha L_h.
\end{equation}

\subsection*{Control of the nonlinear noise contribution}

All identities in this paragraph are on \(\mathcal A_h\).
Put
\[
e=\widetilde f_h-a,\qquad
S_Zb=\divergence(b\nabla Z_h).
\]
Since \(T_{\widehat u_h}=T_v+S_Z\),
\begin{equation}\label{eq:error-equation}
T_{\widehat u_h}e=\mathcal R_h-S_Za.
\end{equation}
The elementary product bounds
\[
\|S_Za\|_2\leq C\|Z_h\|_{H^2},
\qquad
\|S_Za\|_{H^1}
\leq C\big(\|Z_h\|_{H^3}+h^{-1}\|Z_h\|_{H^2}\big)
\]
follow by differentiating \(a\Delta Z_h+\nabla a\cdot\nabla Z_h\)
and using the displayed bounds on \(a\).
Lemma~\ref{lem:transport}, the definition of \(\mathcal A_h\),
and \eqref{eq:residual-bounds} yield the rough estimates
\begin{align}
\|e\|_2
&\leq C\big(B_h+\|Z_h\|_{H^2}\big),
\label{eq:rough-L2}\\
\|e\|_{H^1}
&\leq C\big(
h^{-1}B_h+\|Z_h\|_{H^3}
+h^{-1}\|Z_h\|_{H^2}\big).
\label{eq:rough-H1}
\end{align}

To obtain the sharp variance, rewrite \eqref{eq:error-equation} as
\begin{equation}\label{eq:sharp-expansion}
e=P_v\mathcal R_h-P_vS_Za-P_vS_Ze.
\end{equation}
The middle term is linear in the noise and, by
Lemma~\ref{lem:variance},
\begin{equation}\label{eq:linear-risk}
\E\|P_vS_Za\|_2^2\leq Cn^{-1}h^{-d-3}.
\end{equation}
The last term is controlled without any independence assumption
between the estimated flow and the noise. Namely,
\[
\|P_vS_Ze\|_2
\leq C\left(
\|Z_h\|_{1,\infty}\|e\|_{H^1}
+\|Z_h\|_{2,\infty}\|e\|_2
\right).
\]
Substitute \eqref{eq:rough-L2}--\eqref{eq:rough-H1} and apply
Lemma~\ref{lem:gaussian} and Cauchy--Schwarz. This gives
\begin{equation}\label{eq:nonlinear-risk}
\E\left[\one_{\mathcal A_h}\|P_vS_Ze\|_2^2\right]
\leq C\left(
B_h^2n^{-1}h^{-d-4}L_h
+n^{-2}h^{-2d-8}L_h
\right).
\end{equation}
For clarity, a representative fourth-moment calculation is
\[
\E\left(
\|Z_h\|_{1,\infty}^2\|Z_h\|_{H^3}^2
\right)
\leq
\left(\E\|Z_h\|_{1,\infty}^4\right)^{1/2}
\left(\E\|Z_h\|_{H^3}^4\right)^{1/2}
\leq Cn^{-2}h^{-2d-8}L_h.
\]
The term with
\(h^{-1}\|Z_h\|_{1,\infty}\|Z_h\|_{H^2}\)
has the same bound, as does the term with
\(\|Z_h\|_{2,\infty}\|Z_h\|_{H^2}\).
The terms containing \(B_h\) give the first term in
\eqref{eq:nonlinear-risk}.

Combining \eqref{eq:sharp-expansion}--\eqref{eq:nonlinear-risk},
\[
\E[\one_{\mathcal A_h}\|e\|_2^2]
\leq C\left(
B_h^2+n^{-1}h^{-d-3}
+B_h^2n^{-1}h^{-d-4}L_h
+n^{-2}h^{-2d-8}L_h
\right).
\]
Our bandwidth satisfies \(n^{-1}=h^{2\alpha+d+3}\). Hence the last
display becomes
\[
C\left(
h^{2\alpha}L_h^4+h^{2\alpha}
+h^{2\alpha}L_h^4 h^{2\alpha-1}L_h
+h^{4\alpha-2}L_h
\right)
\leq Ch^{2\alpha}L_h^4.
\]
Here \(\alpha>1\) is used in
\(h^{4\alpha-2}L_h\leq Ch^{2\alpha}\).

Projection onto \([1/2,3/2]\) cannot increase the loss, since the
true conductivity takes values in that interval. By
\eqref{eq:E-stable}, \eqref{eq:completed-bias}, and
\eqref{eq:good-event},
\[
\sup_{f\in\cF(\alpha,r_*)}
\E_f\|\widehat f_n-f\|_2^2
\leq Ch^{2\alpha}L_h^4+Cn^{-3}.
\]
This is the asserted upper bound.

If \(\alpha\notin\mathbb N\), use the versions of the approximation
lemmas without logarithms, so that \(B_h=h^\alpha\).
The logarithms remaining in the stochastic remainder are absorbed
by the strictly positive powers \(h^{2\alpha-1}\) and
\(h^{2\alpha-2}\). The resulting upper bound is \(Ch^{2\alpha}\).
\end{proof}

\section{The lower bound}

We construct pressure perturbations with fine angular oscillations
and fixed radial scale.

\subsection{Angular bumps}

Let \(m=d-1\). Fix a smooth coordinate chart
\(\chi:Q\to\Sph\), and a smaller cube compactly contained in \(Q\).
For sufficiently small \(h\), choose centers \(z_1,\ldots,z_{M_h}\)
in this cube, separated by \(4h\), such that
\[
c h^{-m}\leq M_h\leq C h^{-m}.
\]
Fix a nonzero \(\psi\in C^\infty_c((-1,1)^m)\), and define
\[
\psi_{j,h}(\chi(z))=\psi((z-z_j)/h),
\]
extended by zero outside the chart. Their supports are pairwise
disjoint. Coordinate changes on the fixed chart give
\begin{align}
\|\psi_{j,h}\|_{L^2(\Sph)}^2&\asymp h^m,\label{eq:angular-L2}\\
\|\nabla_{\Sph}\psi_{j,h}\|_{L^2(\Sph)}^2&\asymp h^{m-2}.
\label{eq:angular-grad}
\end{align}
Integration by parts and Cauchy--Schwarz imply
\begin{equation}\label{eq:angular-lap}
\|\Delta_{\Sph}\psi_{j,h}\|_2^2
\geq
\frac{\|\nabla_{\Sph}\psi_{j,h}\|_2^4}
{\|\psi_{j,h}\|_2^2}
\geq c h^{m-4}.
\end{equation}
For \(\theta\in\{0,1\}^{M_h}\), write
\(q_\theta=\sum_j\theta_j\psi_{j,h}\).
Disjointness gives
\[
\|q_\theta\|_2^2\leq C,\qquad
\|\nabla_{\Sph}^kq_\theta\|_\infty\leq C_k h^{-k}
\]
for every fixed \(k\).

\subsection{Realization as conductivities}

Fix a nonnegative, nonzero
\(\rho\in C^\infty_c((1/4,1/2))\).
For a fixed small \(\kappa>0\), define on \(\R^d\)
\[
v_\theta(x)=
\kappa h^{\alpha+2}\rho(|x|)
q_\theta\!\left(\frac{x}{|x|}\right),
\qquad
u_\theta=u_\circ+v_\theta.
\]
The perturbation is zero near \(0\), so this defines a smooth
function everywhere. For each fixed \(k\),
\begin{equation}\label{eq:v-derivatives}
\|\nabla^kv_\theta\|_\infty
\leq C_k\kappa h^{\alpha+2-k}.
\end{equation}
In particular, for sufficiently small \(\kappa\),
\(D^2u_\theta\) is uniformly positive definite, independently of
\(h,\theta\).

\begin{lemma}[Conductivity realization and expansion]\label{lem:realization}
For sufficiently small \(\kappa\), there are smooth positive
conductivities \(f_\theta=1+g_\theta\) such that
\[
\divergence(f_\theta\nabla u_\theta)=1
\quad\text{on }\R^d,
\]
and, uniformly in \(h,\theta\),
\begin{align}
\|\nabla^kg_\theta\|_\infty
&\leq C_k\kappa h^{\alpha-k}
\quad\text{for every fixed }k,\label{eq:g-derivatives}\\
\|g_\theta\|_{C^\alpha(D)}
&\leq C\kappa,\label{eq:g-class}\\
g_\theta&=\ell_\theta+\varepsilon_\theta,\qquad
\ell_\theta=-P_{u_\circ}\Delta v_\theta,\qquad
\|\varepsilon_\theta\|_{L^2(D)}
\leq C\kappa^2h^{2\alpha}.
\label{eq:lower-expansion}
\end{align}
Moreover, \(u_{f_\theta}=u_\theta\) on \(D\).
\end{lemma}

\begin{proof}
The backward-gradient flow of \(u_\theta\) is globally defined on
\(\R^d\): its vector field is globally Lipschitz and has at most
linear growth. Its unique critical point is \(0\), and the flow
contracts towards \(0\).
Define \(f_\theta=P_{u_\theta}\one\) by the same integral as
\eqref{eq:transport-inverse}. That formula and its proof apply on
\(\R^d\) here, since the Hessian has uniform positive bounds and
all derivatives of order at least two are bounded.
Thus \(f_\theta\) is smooth, positive, bounded, and satisfies the
stated equation.

Write \(u=u_\theta\), \(v=v_\theta\), and \(g=f_\theta-1\).
Then
\begin{equation}\label{eq:g-transport}
T_ug=-\Delta v.
\end{equation}
The \(L^\infty\) transport estimate gives
\(\|g\|_\infty\leq C\kappa h^\alpha\).

We prove \eqref{eq:g-derivatives} by induction.
Differentiating \eqref{eq:g-transport} \(k\) times, and collecting
the terms containing \(k\)-th derivatives of \(g\), gives
\[
T_u(\nabla^kg)+\mathscr H_k\nabla^kg=F_k.
\]
Here \(\mathscr H_k\) is the sum of the actions of \(D^2u\) on
the \(k\) tensor indices. If \(D^2u\geq\lambda I\), it is symmetric
and bounded below by \(k\lambda I\). Leibniz's rule gives
\begin{align*}
\|F_k\|_\infty\leq C_k\bigg(
&\|\nabla^{k+2}v\|_\infty\\
&+\sum_{j=2}^k
\|\nabla^{j+1}v\|_\infty
\|\nabla^{k-j+1}g\|_\infty\\
&+\sum_{j=1}^k
\|\nabla^{j+2}v\|_\infty
\|\nabla^{k-j}g\|_\infty
\bigg).
\end{align*}
All derivatives of \(u_\circ\) of order at least three vanish.
Along each trajectory, the homogeneous matrix equation has
propagator norm at most \(e^{-(d+k)\lambda t}\).
Consequently
\[
\|\nabla^kg\|_\infty
\leq ((d+k)\lambda)^{-1}\|F_k\|_\infty.
\]
The differentiability and boundedness needed here also follow by the
variational-equation argument in Lemma~\ref{lem:transport}.
Using the induction hypothesis and \eqref{eq:v-derivatives}, the
right side is bounded by
\[
C_k\big(\kappa h^{\alpha-k}
+\kappa^2h^{2\alpha-k}\big)
\leq C_k'\kappa h^{\alpha-k}.
\]
This proves \eqref{eq:g-derivatives}.

Choose an integer \(J>\alpha\). The dyadic characterization implies
\[
\|g\|_{C^\alpha(\R^d)}
\leq C\left(
\|g\|_\infty+
\|g\|_\infty^{1-\alpha/J}
\|\nabla^Jg\|_\infty^{\alpha/J}
\right).
\]
Indeed, a dyadic block is bounded both by \(C\|g\|_\infty\) and,
using its vanishing moments, by
\(C2^{-jJ}\|\nabla^Jg\|_\infty\); optimize between these bounds.
Substitution of \eqref{eq:g-derivatives} proves
\eqref{eq:g-class}, including at integer \(\alpha\).

Expanding the original equation around \((1,u_\circ)\) gives
\[
T_{u_\circ}g=-\Delta v-\divergence(g\nabla v).
\]
Therefore
\[
g=-P_{u_\circ}\Delta v
  -P_{u_\circ}\divergence(g\nabla v).
\]
The \(L^\infty\) bound for \(P_{u_\circ}\), together with
\eqref{eq:v-derivatives} and \eqref{eq:g-derivatives} at orders
zero and one, yields
\[
\|P_{u_\circ}\divergence(g\nabla v)\|_\infty
\leq C\left(
\|\nabla g\|_\infty\|\nabla v\|_\infty
+\|g\|_\infty\|\Delta v\|_\infty
\right)
\leq C\kappa^2h^{2\alpha}.
\]
This proves \eqref{eq:lower-expansion}.

Finally, \(v_\theta\) vanishes near \(\partial D\), so
\(u_\theta=0\) there. Since \(f_\theta>0\), uniqueness of the
Dirichlet problem implies \(u_{f_\theta}=u_\theta\).
\end{proof}

\subsection{Separation of the conductivities}

For \(x=r\omega\), the unperturbed inverse is explicitly
\[
P_{u_\circ}s(r\omega)
=d r^{-d}\int_0^r t^{d-1}s(t\omega)\,dt.
\]
Using the polar-coordinate Laplacian, for \(r>1/2\) we obtain
\begin{equation}\label{eq:ell-explicit}
\ell_\theta(r\omega)=
-d\kappa h^{\alpha+2}J_\rho r^{-d}
\Delta_{\Sph}q_\theta(\omega),
\qquad
J_\rho=\int_0^\infty t^{d-3}\rho(t)\,dt>0.
\end{equation}
Indeed, the radial part integrates to
\([t^{d-1}\rho'(t)]_0^r=0\).

Let \(H(\theta,\eta)\) be Hamming distance. Since the angular
supports are disjoint, integration over the fixed annulus
\(3/4<|x|<7/8\), using \eqref{eq:angular-lap} and
\eqref{eq:ell-explicit}, gives
\begin{equation}\label{eq:linear-separation}
\|\ell_\theta-\ell_\eta\|_{L^2(D)}^2
\geq c\kappa^2 h^{2\alpha+m}H(\theta,\eta).
\end{equation}

There is a subset \(\Theta_h\subset\{0,1\}^{M_h}\) such that
\begin{equation}\label{eq:code}
\log|\Theta_h|\geq cM_h,\qquad
H(\theta,\eta)\geq M_h/8
\quad(\theta\neq\eta).
\end{equation}
This is the usual binary packing bound; see, for example,
\cite[Lemma~2.9]{Tsybakov}.
One elementary proof is to select words greedily, removing at each
step a Hamming ball of radius \(\lfloor M_h/8\rfloor\).
Such a ball has at most
\[
\sum_{k\leq M_h/8}\binom{M_h}{k}
\leq \exp\{M_h H_b(1/8)\}
\]
points, where \(H_b(t)=-t\log t-(1-t)\log(1-t)<\log2\).
The binomial bound follows by multiplying by
\((t/(1-t))^k\) and applying the binomial theorem.
The greedy procedure therefore leaves exponentially many selected
words and proves \eqref{eq:code}.

Since \(M_h\asymp h^{-m}\), equations
\eqref{eq:linear-separation} and \eqref{eq:code} imply
\[
\|\ell_\theta-\ell_\eta\|_2\geq c\kappa h^\alpha
\quad(\theta,\eta\in\Theta_h,\ \theta\neq\eta).
\]
Lemma~\ref{lem:realization} then gives, for all sufficiently small \(h\),
\begin{equation}\label{eq:nonlinear-separation}
\|f_\theta-f_\eta\|_2\geq c'\kappa h^\alpha.
\end{equation}
The remainder is negligible because
\(\kappa^2h^{2\alpha}=o(\kappa h^\alpha)\).

\subsection{Testing bound}

\begin{proposition}\label{prop:lower}
For every fixed sufficiently small \(r>0\),
\[
R_n(\alpha,r)\geq c\,n^{-2\alpha/(2\alpha+d+3)}
\]
for all sufficiently large \(n\).
\end{proposition}

\begin{proof}
Choose a fixed positive \(\kappa\) so small that all the preceding
constructions apply and \(C\kappa\leq r\), where \(C\) is the constant
in \eqref{eq:g-class}. Thus every \(f_\theta\) belongs to
\(\cF(\alpha,r)\).

Let \(P_\theta\) be the observation law under \(f_\theta\), and let
\(P_\circ\) be that under \(f=1\).
For Gaussian white noise with variance level \(n^{-1}\),
\begin{equation}\label{eq:KL}
\KL(P_\theta,P_\circ)
=\frac n2\|u_\theta-u_\circ\|_2^2
=\frac n2\|v_\theta\|_2^2
\leq Cn\kappa^2h^{2\alpha+4}.
\end{equation}
The equality follows directly from the Gaussian likelihood
\[
\frac{dP_\theta}{dP_\circ}
=
\exp\left\{
n\big(Y_n(v_\theta)-\langle u_\circ,v_\theta\rangle\big)
-\frac n2\|v_\theta\|_2^2
\right\}.
\]
It can equivalently be obtained by finite-dimensional projections
and convergence of the isonormal series. The inequality in
\eqref{eq:KL} uses \(\|q_\theta\|_2^2\leq C\).

Take
\[
h=n^{-1/(2\alpha+d+3)}.
\]
Since \(m=d-1\),
\[
nh^{2\alpha+4}=h^{-m},
\qquad
\log|\Theta_h|\geq c h^{-m}.
\]
Decrease the fixed \(\kappa\), if necessary, so that
\begin{equation}\label{eq:information-small}
\frac1{|\Theta_h|}
\sum_{\theta\in\Theta_h}\KL(P_\theta,P_\circ)
\leq\frac18\log|\Theta_h|.
\end{equation}

We recall the finite-hypothesis information inequality used here,
a form of Fano's inequality; compare
\cite[Theorem~2.5 and Chapter~2]{Tsybakov}.
Let \(J\) be uniform over \(N\) hypotheses and let
\(\overline P=N^{-1}\sum_jP_j\).
For every reference law \(P_0\),
\[
I(J;Y)
=\frac1N\sum_j\KL(P_j,\overline P)
\leq\frac1N\sum_j\KL(P_j,P_0),
\]
because the difference of the right side and the middle expression
is \(\KL(\overline P,P_0)\).
For any classifier \(\widehat J\), with error probability \(p_e\),
the binary error indicator gives
\[
H(J\mid Y)\leq\log2+p_e\log N.
\]
Since \(H(J)=\log N\), it follows that
\[
p_e\geq1-\frac{I(J;Y)+\log2}{\log N}.
\]
Applying this with \eqref{eq:information-small} shows that every
classifier has average error probability at least \(1/2\) for
large \(n\).

Given any \(L^2(D)\)-valued estimator, classify by its nearest member
of \(\{f_\theta:\theta\in\Theta_h\}\), resolving ties by a fixed rule.
By \eqref{eq:nonlinear-separation}, misclassification implies an
\(L^2\)-error of at least \(c\kappa h^\alpha/2\).
Consequently,
\[
\sup_{f\in\cF(\alpha,r)}
\E_f\|\widehat f-f\|_2^2
\geq
\frac1{|\Theta_h|}
\sum_{\theta\in\Theta_h}
\E_\theta\|\widehat f-f_\theta\|_2^2
\geq c\kappa^2h^{2\alpha}.
\]
Taking the infimum over estimators proves the proposition.
\end{proof}

\begin{proof}[Proof of Theorem~\ref{thm:main}]
Choose \(r_*\) sufficiently small for
Lemma~\ref{lem:forward} and Proposition~\ref{prop:upper}.
The upper bound then holds for every subclass
\(\cF(\alpha,r)\), \(0<r\leq r_*\).
Proposition~\ref{prop:lower}, with its fixed \(\kappa\) chosen
according to \(r\), gives the matching lower bound.
\end{proof}

\begin{remark}
The effective denominator is \(2\alpha+d+3\), rather than the
denominator obtained by treating the forward map as an isotropic
operator of a single differential order. In the lower bound there
are \(d-1\) angular directions and a two-derivative attenuation.
In the upper bound, transport integration reduces the integrated
variance of a twice differentiated, smoothed white noise from
order \(n^{-1}h^{-d-4}\) to order \(n^{-1}h^{-d-3}\).
The small-radius assumption supplies uniform convexity of the
pressure and uniform control of this transport geometry.
\end{remark}

\begin{thebibliography}{9}

\bibitem{GT}
D.~Gilbarg and N.~S.~Trudinger,
\emph{Elliptic Partial Differential Equations of Second Order},
second edition, Classics in Mathematics,
Springer, Berlin, 2001.
doi:10.1007/978-3-642-61798-0.

\bibitem{Rychkov}
V.~S.~Rychkov,
On restrictions and extensions of the Besov and Triebel--Lizorkin
spaces with respect to Lipschitz domains,
\emph{Journal of the London Mathematical Society} (2)
\textbf{60} (1999), 237--257.
doi:10.1112/S0024610799007723.

\bibitem{Taylor}
M.~E.~Taylor,
\emph{Partial Differential Equations III: Nonlinear Equations},
second edition, Applied Mathematical Sciences, vol.~117,
Springer, New York, 2011.

\bibitem{Tsybakov}
A.~B.~Tsybakov,
\emph{Introduction to Nonparametric Estimation},
Springer Series in Statistics,
Springer, New York, 2009.
doi:10.1007/b13794.

\end{thebibliography}

\end{document}
